\section*{Hoofdstuk \ref{ch:b2:global-change} --- Mondiale verandering en de biosfeer}

\begin{solution}{exo:b2:global-change:1}
$\Delta T = \lambda C_{\text{cum}}$ met $\lambda = \qty{1.65}{\celsius}$
per \qty{1000}{GtC}: \qty{1.2}{\celsius}, \qty{1.65}{\celsius},
\qty{3.3}{\celsius}.
\end{solution}

\begin{solution}{exo:b2:global-change:2}
Graaddagen: de som over de dagen van het overschot van de daggemiddelde
temperatuur boven een basis, de \omterm{thm:b2:global-change:phenology}{thermische tijd} die de ontwikkeling
van planten en ectotherme dieren bepaalt. Ontkoppeling: wanneer twee
interagerende soorten hun seizoenen op verschillende signalen afstemmen en
opwarming de ene verschuift en de andere niet --- de vliegenvanger, die zich in
Afrika op de daglengte richt, komt op zijn oude datum aan en vindt de
rupsenpiek, die zich op de temperatuur richt, al twee weken voorbij.
\end{solution}

\begin{solution}{exo:b2:global-change:3}
Ongeveer \qty{150}{m} bergop en \qty{150}{km} naar de pool per graad. Een
soort op de top heeft geen hoger gelegen grond: haar isotherm stijgt
boven de berg uit en haar habitat verdwijnt, terwijl het oppervlak van
elke band onderweg naar boven met de hoogte krimpt.
\end{solution}

\begin{solution}{exo:b2:global-change:4}
Opgelost \ensuremath{\mathrm{CO_2}} vormt koolzuur, dat
waterstofionen vrijmaakt, en de pH daalt evenredig met de
logaritme van de \ensuremath{\mathrm{CO_2}}-druk. De extra
waterstofionen zetten carbonaat om in bicarbonaat, zodat de
carbonaatconcentratie en de \omterm{thm:b2:global-change:acid}{verzadigingstoestand} van calciumcarbonaat dalen.
Organismen die aragoniet opbouwen --- koralen, pteropoden, veel larven
--- en calciet --- coccolithoforen, foraminiferen, weekdieren ---
moeten meer investeren om schelpen te bouwen en verliezen ze onder de verzadiging.
\end{solution}

\begin{solution}{exo:b2:global-change:5}
\qty{1.5}{\celsius}: $1.5/1.65\times 1000 = \qty{909}{GtC}$ in totaal,
\qty{209}{GtC} resterend, 21 jaar. \qty{2}{\celsius}:
\qty{1212}{GtC}, \qty{512}{GtC} resterend, 51 jaar.
\end{solution}

\begin{solution}{exo:b2:global-change:6}
$\sqrt{2\times 150/0.12} = 50$ dagen; vervroeging $1.5/0.12 = 12.5$
dagen.
\end{solution}

\begin{solution}{exo:b2:global-change:7}
\qty{3}{\celsius}: de isothermen stijgen $3000/6.5 = \qty{460}{m}$; de
band verschuift naar \qtyrange{1260}{1860}{m}, maar de berg eindigt op
\qty{1800}{m}: de bovenste \qty{60}{m} van de band gaan verloren en de rest
ligt waar het oppervlak $2^{-460/300} = 0.35$ is van wat het was --- ongeveer
twee derde van het oppervlak verdwenen. \qty{5}{\celsius}: \qty{770}{m}; de
band zou \qtyrange{1570}{2170}{m} zijn; alleen \qtyrange{1570}{1800}{m}
bestaat, een derde van de breedte van de band, op $2^{-770/300} = 0.17$ van het
oppervlak --- minder dan een tiende van het oorspronkelijke oppervlak: de soort
is in feite verdwenen.
\end{solution}

\begin{solution}{exo:b2:global-change:8}
560: pH $7.93$, waterstofionen $\times 1.9$; 800: $7.79$, $\times
2.6$; 1000: $7.70$, $\times 3.1$.
\end{solution}

\begin{solution}{exo:b2:global-change:9}
$1 - (1 - h)^{0.25}$: \qty{26}{\%}, \qty{53}{\%}, \qty{82}{\%}. Het
onmiddellijke verlies is kleiner omdat het restant nog populaties van de
meeste soorten bevat; die zijn te klein om stand te houden en
sterven over decennia uit --- de \omterm{thm:b2:global-change:sar}{extinctieschuld}.
\end{solution}

\begin{solution}{exo:b2:global-change:10}
$c = 2\sqrt{0.7\times 20} = \qty{7.5}{km/yr}$; \qty{1000}{km} in
ongeveer 130 jaar. Het halveren van $r$ of van $D$ deelt $c$ door
$\sqrt 2$, tot \qty{5.3}{km/yr}; geen van beide halveert de snelheid --- om de
snelheid te halveren moet men het product $rD$ door vier delen.
\end{solution}

\begin{solution}{exo:b2:global-change:11}
Achtergrond: $8\times 10^{6}\times 10^{-6} = 8$ per jaar, 800 per
eeuw; bij duizend keer 8000 per jaar, \num{800000} per eeuw.
Tienduizend gedocumenteerde in een eeuw is ongeveer tien keer de
achtergrond, geen duizend keer --- maar de documentatie dekt alleen de
enkele procenten soorten die beschreven en gevolgd worden (vogels,
zoogdieren, enkele planten), en daaronder is de snelheid inderdaad honderd
keer de achtergrond of meer; voor de onbeschreven meerderheid worden de meeste
uitstervingen nooit gezien, en is de schatting via de \omterm{thm:b2:global-change:sar}{soorten-oppervlakterelatie}
het enige houvast. De waarheid ligt tussen het getelde en het
geschatte, en daarom is het opgegeven bereik zo breed.
\end{solution}

\begin{solution}{exo:b2:global-change:12}
De temperatuur is evenredig met de \omterm{prop:b2:global-change:tcre}{cumulatieve uitstoot}; wanneer de
netto-uitstoot nul is, stopt het cumulatieve totaal met groeien, en de
opwarming ook. In het \omterm{thm:b2:biogeochemical-cycles:box}{boxmodel} begint het atmosferische overschot bij een
uitstoot van nul te dalen doordat de putten blijven werken, wat zou
afkoelen --- maar de oceaan, die nog naar het evenwicht met het verhoogde
\ensuremath{\mathrm{CO_2}} opwarmt, blijft het oppervlak verwarmen, en de twee
effecten heffen elkaar bijna op: de temperatuur blijft eeuwenlang
waar ze is. Wat aan het argument ontbreekt: de andere broeikasgassen
(door de korte levensduur van methaan daalt zijn opwarming snel zodra de
uitstoot stopt), de aerosolen die nu een deel van de opwarming maskeren en
binnen enkele weken verdwijnen na de uitstoot die ze maakt, en de
terugkoppelingen --- koolstof uit permafrost, afsterven van bossen --- die zelf
uitstoot zouden kunnen toevoegen.
\end{solution}

\begin{solution}{pb:b2:global-change:1}
\textbf{1.} A: $700 + 80\times 10 = \qty{1500}{GtC}$. B: $700 +
\tfrac12\times 50\times 10 = \qty{950}{GtC}$.
\textbf{2.} A: $1.65\times 1.5 = \qty{2.5}{\celsius}$. B:
$1.65\times 0.95 = \qty{1.6}{\celsius}$.
\textbf{3.} 2050. A: \qty{1000}{GtC}, \qty{1.65}{\celsius}. B:
uitstoot $10\int_0^{30}(1 - t/50)\,\mathrm{d}t = 10(30 - 9) =
\qty{210}{GtC}$, totaal 910, \qty{1.5}{\celsius}.
\textbf{4.} $\lambda\times 100 = \qty{0.165}{\celsius}$ per decennium,
iets onder de waargenomen $0.2$ (die de andere gassen en het wegvallen van de
maskering door aerosolen omvat).
\textbf{5.} Ja: de \omterm{prop:b2:global-change:tcre}{cumulatieve uitstoot} stopt in 2070 met groeien, en de
opwarming ook, die op \qty{1.6}{\celsius} blijft --- de
evenredigheid betekent dat de temperatuur door het totaal wordt bepaald, niet
door het tempo.
\textbf{6.} $2/1.65\times 1000 = \qty{1212}{GtC}$; pad A bereikt die
na $512/10 = 51$ jaar, in 2071.
\textbf{7.} $\sqrt{2\times 200/0.1} = 63$ dagen.
\textbf{8.} Opwarming ten opzichte van 2020: A $2.5 - 1.2 =
\qty{1.3}{\celsius}$, B \qty{0.4}{\celsius}. Vervroeging $\Delta T/b$:
A 13 dagen, B 4 dagen.
\textbf{9.} Kloof: A $20 + 13 = 33$ dagen; B 24 dagen.
\textbf{10.} Rupsen beschikbaar van dag 15 tot dag 25 na de
bladontplooiing. B: de vogel komt aan op dag 24, binnen het venster, nipt.
A: dag 33, acht dagen na de laatste rups.
\textbf{11.} De vogel vervroegt $3\times 1.3 = 4$ dagen: kloof $33 - 4 =
29$ dagen --- nog steeds buiten het venster.
\textbf{12.} De vogel kan bij de nieuwe temperatuur leven; wat hij
niet kan, is kuikens voeren wanneer het voedsel verdwenen is, omdat zijn signaal
en dat van zijn prooi verschillend op opwarming reageren.
\textbf{13.} A: $1.3/6.5\times 1000 = \qty{200}{m}$ omhoog en
$1.3/0.65\times 100 = \qty{200}{km}$ naar de pool. B: \qty{57}{m} en
\qty{57}{km}.
\textbf{14.} In 80 jaar verplaatst de boom zich \qty{8}{km}: hij kan
\qty{200}{km} (A) niet volgen en schiet ook voor \qty{57}{km} (B) tekort ---
bomen lopen achter, en hun arealen zullen op beide paden eeuwenlang niet in
evenwicht met het klimaat zijn.
\textbf{15.} A: de band zou \qtyrange{1700}{2000}{m} zijn: hij
bereikt net de top en kan nergens meer heen. De berg versmalt naar boven,
zodat dezelfde band nog maar $2^{-200/300} = 0.63$ keer zijn vroegere
oppervlak beslaat --- en verdere opwarming heeft geen grond meer om hem
naartoe te verschuiven.
\textbf{16.} B: \qty{57}{m} omhoog, oppervlaktefactor $2^{-57/300} = 0.88$: een
verlies van \qty{12}{\%}.
\textbf{17.} Acht decennia met \qty{30}{km}: \qty{240}{km}, meer dan
de \qty{200}{km} die op A nodig is: de vis houdt gelijke tred (en verlaat
zijn oude visgrond).
\textbf{18.} Aan de ondergrens ontmoet de soort concurrenten,
ziekteverwekkers en predatoren die van beneden opschuiven en die ze voordien
niet tegenkwam, en haar eigen fysiologie, aangepast aan de oude temperatuur,
wordt door die van hen overtroffen: de terugtrekking is evenzeer competitief als
fysiologisch.
\textbf{19.} A: $\mathrm{pH} = 8.2 - 0.9\log_{10}(600/280) = 7.90$;
waterstofionen $10^{0.30} = 2.0$ keer die van 1750. B: $8.04$; $1.44$
keer.
\textbf{20.} $\Omega = 4/(\text{verhouding})$: A $2.0$, B $2.8$. Bij 2
verkalken de koralen nog, maar traag, met zwakkere skeletten en riffen die
sneller eroderen dan ze groeien; bij $2.8$ bouwen ze, al verbleekt de warmte
van pad B ze nog steeds.
\textbf{21.} A: $400\times 0.4^{0.25} = 318$, 82 verloren. B:
$400\times 0.8^{0.25} = 378$, 22 verloren.
\textbf{22.} A: $318\times(1 - 0.05\times 1.3) = 298$. B:
$378\times(1 - 0.05\times 0.4) = 371$.
\textbf{23.} A: $1 - 298/400 = \qty{26}{\%}$. B: \qty{7}{\%}.
\textbf{24.} Op beide paden de ontbossing: 82 tegen 20 soorten op A, 22
tegen 7 op B --- de opwarmingsterm van het model is bescheiden; in werkelijkheid
werken de twee op elkaar in, omdat een versnipperd bos zijn areaal niet kan
verschuiven.
\textbf{25.} Pad A: \qty{2.5}{\celsius}, pH $7.90$, een kwart van de
boomsoorten verloren. Pad B: \qty{1.6}{\celsius}, pH $8.04$, \qty{7}{\%}
verloren.
\end{solution}
